
\documentclass{article}
\usepackage{colortbl}
\usepackage{makecell}
\usepackage{multirow}
\usepackage{supertabular}

\begin{document}

\newcounter{utterance}

\centering \large Interaction Transcript for game `cladder', experiment `full\_v1.5\_default', episode 2103 with selfrefine{-}Qwen3.5{-}27B.
\vspace{24pt}

{ \footnotesize  \setcounter{utterance}{1}
\setlength{\tabcolsep}{0pt}
\begin{supertabular}{c@{$\;$}|p{.15\linewidth}@{}p{.15\linewidth}p{.15\linewidth}p{.15\linewidth}p{.15\linewidth}p{.15\linewidth}}
    \# & \multicolumn{2}{c}{Player} && \multicolumn{2}{c}{Game Master} \\
    \hline

    \theutterance \stepcounter{utterance}  
    & & & \multicolumn{4}{p{0.6\linewidth}}{
        \cellcolor[rgb]{0.9,0.9,0.9}{
            \makecell[{{p{\linewidth}}}]{
                \texttt{\tiny{[P1$\langle$GM]}}
                \texttt{You are an expert in causal inference. The following question is not a typical commonsense query, but rather a meticulously designed question created by a professor specializing in causal inference, intended to assess the students' mastery of the course content.} \\
\texttt{Imagine a self{-}contained, hypothetical world with only the following conditions, and without any unmentioned factors or causal relationships: Gender has a direct effect on department competitiveness and admission status. Department competitiveness has a direct effect on admission status. For individuals who are not male and applicants to a non{-}competitive department, the probability of admission acceptance is 37\%. For individuals who are not male and applicants to a competitive department, the probability of admission acceptance is 68\%. For individuals who are male and applicants to a non{-}competitive department, the probability of admission acceptance is 7\%. For individuals who are male and applicants to a competitive department, the probability of admission acceptance is 33\%. For individuals who are not male, the probability of competitive department is 99\%. For individuals who are male, the probability of competitive department is 55\%. Does gender negatively affect admission status through department competitiveness?} \\
\texttt{Start your answer with 'Yes' or 'No', followed by additional reasoning or evidence to support your explanation.} \\
            }
        }
    }
     \\ \\

    \theutterance \stepcounter{utterance}  
    & \multicolumn{4}{p{0.6\linewidth}}{
        \cellcolor[rgb]{0.561,0.737,0.561}{
            \makecell[{{p{\linewidth}}}]{
                \texttt{\tiny{[P1$\rangle$GM]}}
                \texttt{No} \\
\\ 
\texttt{To determine whether gender negatively affects admission status through department competitiveness, we must calculate the **Indirect Effect** (or Mediated Effect) of Gender on Admission via Department Competitiveness. This effect is the product of the effect of Gender on the mediator (Department Competitiveness) and the effect of the mediator on the outcome (Admission).} \\
\\ 
\texttt{Let's define the variables:} \\
\texttt{*   \$G\$: Gender (\$1\$ = Male, \$0\$ = Not Male)} \\
\texttt{*   \$D\$: Department Competitiveness (\$1\$ = Competitive, \$0\$ = Non{-}competitive)} \\
\texttt{*   \$Y\$: Admission Status (\$1\$ = Accepted, \$0\$ = Rejected)} \\
\\ 
\texttt{**1. Analyze the Effect of Gender on Department Competitiveness (\$G \textbackslash{}to D\$)**} \\
\texttt{We compare the probability of applying to a competitive department for Males versus Not Males.} \\
\texttt{*   \$P(D=1 | G=0) = 0.99\$ (Not Males are 99\% likely to apply to Competitive departments).} \\
\texttt{*   \$P(D=1 | G=1) = 0.55\$ (Males are 55\% likely to apply to Competitive departments).} \\
\texttt{*   **Effect:** Being Male (\$G=1\$) **decreases** the probability of applying to a competitive department by \$0.55 {-} 0.99 = {-}0.44\$.} \\
\\ 
\texttt{**2. Analyze the Effect of Department Competitiveness on Admission (\$D \textbackslash{}to Y\$)**} \\
\texttt{We examine how applying to a competitive department affects admission rates, holding gender constant.} \\
\texttt{*   **For Not Male (\$G=0\$):**} \\
\texttt{    *   \$P(Y=1 | G=0, D=1) = 0.68\$} \\
\texttt{    *   \$P(Y=1 | G=0, D=0) = 0.37\$} \\
\texttt{    *   Effect: \$0.68 {-} 0.37 = +0.31\$. (Competitive departments increase admission chances).} \\
\texttt{*   **For Male (\$G=1\$):**} \\
\texttt{    *   \$P(Y=1 | G=1, D=1) = 0.33\$} \\
\texttt{    *   \$P(Y=1 | G=1, D=0) = 0.07\$} \\
\texttt{    *   Effect: \$0.33 {-} 0.07 = +0.26\$. (Competitive departments increase admission chances).} \\
\\ 
\texttt{In this specific hypothetical world, applying to a "competitive" department is a **positive** factor for admission for both genders.} \\
\\ 
\texttt{**3. Calculate the Indirect Effect**} \\
\texttt{The indirect effect of being Male on Admission through Department Competitiveness is the combination of:} \\
\texttt{1.  Being Male reduces the likelihood of applying to a competitive department (} \\
            }
        }
    }
    & & \\ \\

    \theutterance \stepcounter{utterance}  
    & & & \multicolumn{2}{p{0.3\linewidth}}{
        \cellcolor[rgb]{0.9,0.9,0.9}{
            \makecell[{{p{\linewidth}}}]{
                \texttt{\tiny{[GM$|$GM]}}
                \texttt{No} \\
            }
        }
    }
    & & \\ \\

    \theutterance \stepcounter{utterance}  
    & & & \multicolumn{2}{p{0.3\linewidth}}{
        \cellcolor[rgb]{0.9,0.9,0.9}{
            \makecell[{{p{\linewidth}}}]{
                \texttt{\tiny{[GM$|$GM]}}
                \texttt{yes} \\
            }
        }
    }
    & & \\ \\

    \theutterance \stepcounter{utterance}  
    & & & \multicolumn{2}{p{0.3\linewidth}}{
        \cellcolor[rgb]{0.9,0.9,0.9}{
            \makecell[{{p{\linewidth}}}]{
                \texttt{\tiny{[GM$|$GM]}}
                \texttt{game\_result = LOSE} \\
            }
        }
    }
    & & \\ \\

\end{supertabular}
}

\end{document}
